\chapter{Discussion and Threats to Validity}\label{ch:discussion}

\section{What the three results together establish}\label{sec:disc-arc}

Taken singly, each of the preceding results answers a bounded question; taken
together they trace a single arc, from inference through composition to
integration. \Cref{ch:inferrer} establishes that specifications can be recovered
cheaply---the JML-Inferrer recovers sound, verifiable contracts from unannotated
Java by syntax-driven heuristics, and supplying those contracts to a language
model measurably improves the unit tests it generates. \Cref{ch:compositional}
establishes that the same inferred contracts can be reasoned \emph{about}: a
compositional, weakest-precondition pass over the call graph strictly extends the
local heuristics, and---run differentially across two versions of a dependency's
contract---turns behavioural version-compatibility into a static, run-free
analysis. \Cref{ch:integration} shows what those capabilities make possible once
assembled: a reference architecture in which a machine-checkable contract is the
load-bearing pivot between ambiguous human intent and precise machine
verification, serving at once as the input an agent builds against and the oracle
its output is checked against. The movement is from cheap specifications, to
reasoning about compatibility, to specifications made load-bearing in a lifecycle.

The three research questions may now be assessed against that arc. \textbf{RQ1}
asked how faithfully specifications can be inferred from unannotated code and
whether the result is useful downstream; it is answered by \cref{ch:inferrer} and
answered by \emph{measurement}. Fidelity is anchored to formal verification rather
than inspection---where the engine emits a contract and OpenJML~\citep{cok2014openjml}
discharges it through Z3~\citep{demoura2008z3}, the contract is sound with respect
to the code---a stronger guarantee than the likely-invariant detection of dynamic
tools such as Daikon~\citep{ernst2007daikon}, and the downstream benefit is
established by a controlled comparison rather than asserted. \textbf{RQ2} asked
whether a compositional pass extends inference and whether it enables
version-compatibility reasoning; it is answered by \cref{ch:compositional}, in part
by measurement (the additional well-formed preconditions recovered on a large
corpus) and in part by a design validated by instantiation (the differential
blast-radius analysis). The link is not analogy but identity: the arrow that, run
forwards within a version, strengthens a caller's specification is, run across
versions, the edge that carries a breaking change from a library to its clients.
\textbf{RQ3} asked whether inferred contracts can serve as the contract layer and
verification loop of AI-native lifecycles; it is addressed by \cref{ch:integration}
as a reference architecture that is \emph{designed, not deployed}. The architecture
is concrete and its components are individually evidenced by the earlier chapters,
but the end-to-end lifecycle is demonstrated for feasibility and coherence, not
established by a controlled effect size. That distinction is maintained
deliberately throughout, and the industrial pilot that would supply such a size is
future work (\cref{ch:conclusion}).

\section{Threats to validity}\label{sec:disc-threats}

The value of a candid thesis lies partly in its restraint, and several threats
bound the claims above.

\paragraph{Corpus generalisability (external).} The quantitative evaluation, and
the compositional-reach measurement in particular, leans on Apache Commons Lang and
similarly structured Java libraries. These are real, widely-depended-upon codebases
rather than synthetic benchmarks, so the results are representative of an important
class of production Java; they are \emph{not} random samples of all Java, and the
thesis does not claim that the reported figures transfer unchanged to arbitrary
corpora. A library written in a markedly different idiom---heavy reflection, dynamic
proxies, concurrency-dominated control flow---may expose patterns the heuristics do
not recognise and obligations the solver cannot discharge. Multi-corpus replication
is the obvious next step.

\paragraph{Heuristic partiality (internal).} The inference engine recognises the
patterns it was built to recognise and is silent elsewhere, so its inferred
contracts are sound \emph{fragments}, not complete specifications. This biases the
population of contracts toward the tractable and could inflate apparent fidelity if
the easy cases are also the verifiable ones. The mitigation is to report fidelity as
a rate conditioned on emission and verification, and to record the engine's silence
honestly rather than count it as success: the thesis claims soundness where a
contract is emitted and verified, never coverage of all behaviours.

\paragraph{SMT boundaries (conclusion).} The verification spine rests on Z3, and
the underlying logic is undecidable. The programme repeatedly encounters timeouts on
multiplicative arithmetic, quantified-array-theory queries, and rich preconditions;
where the solver cannot decide, the pipeline neither confirms nor refutes. An honest
account records such obligations as \emph{unverified}, not as results the tooling
did not deliver. Operationally the configuration is tuned to maximise the fraction
discharged within budget, and the fail-open, sound-in-alarm posture ensures that an
undecided obligation degrades to ``unverified'' or ``unflagged'' rather than to a
false guarantee. For the compatibility analysis this makes the reading asymmetric
and must be stated as such: an alarm is trustworthy, but silence is not a guarantee
of compatibility, because a break may be missed when a callee contract is inferred
too weakly to capture the depended-upon behaviour, or when a timeout leaves the
relevant obligation unproved.

\paragraph{Verified versus well-formed (construct).} The compositional preconditions
recovered in \cref{ch:compositional}---roughly $18{,}854$ on the corpus
studied---are \emph{well-formed}: syntactically
valid and scope-correct, admitted only when their entry-state scope is satisfiable.
They are \emph{not} separately discharged by extended static checking, and the thesis
never blurs the two. A verified specification is one the checker has proved against
the code; a well-formed one has passed a weaker, structural filter. Reporting the
well-formedness-filtered count as the headline figure is itself part of the
mitigation, since it excludes vacuous contracts---an inference emitting
\texttt{ensures true} is trivially verifiable yet says nothing---from being counted
as wins.

\paragraph{Construct validity of the utility measure.} The downstream-utility study
of \cref{ch:inferrer} measures test quality by mutation score, a proxy for the
oracle strength of a generated suite~\citep{barr2015oracle} rather than a direct
measurement of correctness caught in the field. A proxy can diverge from the
construct it stands for. The mitigation is that mutation score is an \emph{external}
effect a vacuous contract cannot produce, giving a check on construct validity
independent of the contract's internal form. A related internal threat is
run-to-run variance in the language model: its output is stochastic, so a
test-generation result depends on a model whose behaviour varies between runs. This
is mitigated by the paired design---the same prompts with and without the contract
treatment, so non-determinism is shared across conditions and the treatment effect
is isolated---but the study uses a single model family, so it establishes direction
and plausibility rather than a universal effect size.

\paragraph{Composed validity.} Finally, the validity of a composed system is not the
conjunction of the validities of its parts. Each chapter's result holds under its own
assumptions, and those assumptions are not identical; an assumption satisfied in
isolation may be violated in combination. The reference lifecycle of
\cref{ch:integration} is accordingly claimed as a design coherent in its parts, with
the interaction effects between them identified as precisely the questions a future
end-to-end evaluation must measure. Naming that gap is part of the contribution: it
tells a future evaluator where to look.
